% Source printed pp. iii-v, PDF pp. 3-5. Complete preface.
\chapter*{前言}
\addcontentsline{toc}{chapter}{前言}
\label{preface}
这些讲义建立代数几何中模理论的基础, 目标是在特征 $0$ 的情形下, 自足地证明以下定理: 

\phantomsection\noindent\textbf{定理 A.}\label{thm:A}\quad 亏格 $g\geq 2$ 的稳定曲线的模叠 $\overline{\mathcal M}_g$ 是一个光滑, 固有且不可约的 Deligne--Mumford 叠, 维数为 $3g-3$, 并且具有射影的粗模空间 $\overline M_g$.

\medskip
\phantomsection\noindent\textbf{定理 B.}\label{thm:B}\quad 设 $C$ 是一条亏格为 $g$ 的光滑, 连通射影曲线. $C$ 上秩为 $r$, 次数为 $d$ 的半稳定向量丛的模叠 $\cBun_{r,d}^{\mathrm{ss}}(C)$ 是一个光滑, 普遍闭且不可约的代数叠, 维数为 $r^2(g-1)$, 并且具有射影的良模空间 $M_{r,d}^{\mathrm{ps}}(C)$.

在此过程中, 我们建立代数空间与叠的基础, 为讨论和确立模空间的几何性质提供便利的语言. 引入这些基础需要同时展开若干主题, 包括: 
\begin{itemize}
\item 采用代数几何中的函子观点: 概形 $X$ 由它同所有其他概形的关系决定, 也就是由其可表函子 $\Sch\to\Sets$, $T\mapsto\Mor(T,X)$ 决定; 
\item 把叠引入为编码对象对称性的范畴工具, 例如, 把 $\mathcal M_g$ 看成这样的对应: 给每个概形 $T$ 赋予 $T$ 上所有光滑曲线族构成的群胚; 
\item 用平展拓扑代替概形上的 Zariski 拓扑: 我们将引入 Grothendieck 拓扑作为拓扑空间的推广, 并系统地使用平展态射的下降理论; 
\item 借助若干通常不会在第一门代数几何课程中涉及的进阶主题, 例如概形的平坦, 平展和光滑态射的性质, 代数群及其作用, 形变理论, 以及曲面的双有理几何.
\end{itemize}
为这些基础内容选择一个线性的叙述次序并非易事. 为减轻这一困难, 我们把大量背景知识放在附录中, 使讲义正文始终围绕上述两个定理, 集中讨论模理论.

\paragraph*{预备知识.} 我们假定你熟悉概形, 例如达到 \cite[第 $2$--$3$ 章]{Har77} 或 \cite[第 II--V 部分]{Vak17} 的程度, 但不要求更多. 另一方面, 你会在模理论中看到许多重要的代数几何概念如何实际发挥作用, 例如上同调与基变换, 固有性的赋值判别法, 以及形变理论. 这也提供了补足背景知识, 加深对概形理论核心概念理解的机会. 事实上, \cite{Har77,Vak17} 的每一节都会同本书至少某一部分有关.

\paragraph*{来源.} 模空间有着丰富的历史, 其现代起源可追溯至 Bernhard Riemann 在 $19$ 世纪的工作(见第~\bookref{1.1}~节). 本书发展模理论的方式, 源自 Alexander Grothendieck, David Mumford, Michael Artin, János Kollár, Johan de Jong 以及许多其他人的工作. 本书许多主题的主要资料来源是 Stacks Project~\cite{SP}. 稳定曲线模空间的研究发端于 \cite{GIT,DM69}, 而教材 \cite{HM98,ACG11} 给出了进一步的细节. 类似地, 半稳定向量丛模空间的研究始于 \cite{Mum63,Ses67}, 书籍 \cite{LP97,HL10} 则提供了更详尽的论述.

\paragraph*{我们的路线.} 我们选择一条通往 $\overline M_g$ 山峰(定理~\hyperref[thm:A]{A})和 $\Bun_C$ 山峰(定理~\hyperref[thm:B]{B})的缓坡路线, 尽量减少技术性的陡壁和冰隙. 沿途可以不断眺望前方的山顶, 以及周围的山谷和峰峦. 附录提供攀登所需技术装备的基础训练, 我们还会先热身攀登 Hilbert 与 Quot 两座山峰(定理~\bookref{2.1.2}~和~\bookref{2.1.3}), 以熟悉这些装备. 尽管登山时极不建议这样做, 你也可以不带任何装备, 不经过训练便直接奔向山顶, 等到跌入冰隙时再去学习所需的背景知识.

本书秉持这样一种理念: 对数学的理解, 最好通过深入研究例子来获得. 因此, 我们通过三个主要例子说明模理论的概念: Hilbert 与 Quot 概形(没有自同构的对象的模空间), 稳定曲线的模叠 $\overline{\mathcal M}_g$(具有有限自同构), 以及固定曲线上半稳定向量丛的模叠 $\cBun_{r,d}^{\mathrm{ss}}(C)$(具有无限自同构). 在研究高维簇的模, 或曲面上 Bridgeland 半稳定复形的模之前, 当然应当扎实地理解这些例子.

本书很长(抱歉!), 不应从头到尾按线性次序阅读. 我们希望每一章, 乃至每一节, 都能够独立阅读, 读者可以按自己的意愿跳读. 全书力图为每一个新概念提供例子, 反例和习题, 并总是从最初等的情形开始. 习题的类型各不相同, 例如计算更多例子, 证明进一步的性质, 或补全证明中的技术细节, 但目的始终相同: 让读者有机会练习并尝试操作这些新的数学对象.

Stacks Project 的做法是以最大的普遍性证明定理, 而我们常常会牺牲一定程度的普遍性, 以换取更加直观或几何的证明. 最显著的是, 虽然我们尽可能在任意域 $\Bbbk$(或基概形 $S$)上工作, 但在曲线的稳定约化定理~\bookeqref{6.5.1}~的证明, 以及半稳定向量丛模空间的构造中, 都假定域 $\Bbbk$ 的特征为 $0$. 有时我们会加入并非严格必要的 Noether 假设, 例如代数叠的固有性赋值判别法~\bookeqref{4.8.7}; 或者只对 Deligne--Mumford 叠, 而非具有有限惯性的代数叠, 建立某个定理, 例如 Keel--Mori 定理~\bookeqref{5.4.6}.

只要可能, 我们就给出指向 Stacks Project 及原始文献的指引, 那里包含更一般的结果或其他论证. 我们希望本书能成为阅读 Stacks Project 中更加详尽论述的指南. 通过本书介绍的工具和技巧, 读者应当有能力研究这些模空间的几何性质, 如 \cite{HM98,HL10} 中所做的那样.

\paragraph*{致谢.} 本书源于作者在 $2021$ 年冬季学期于华盛顿大学讲授的线上研究生课程 Math 582C: 叠与模空间导论. 此后几年间, 我有机会在华盛顿大学讲授类似的研究生课程, 并在埃塞克斯大学, 智利天主教大学, 比萨高等师范学校, 杜伊斯堡--埃森大学, 勃艮第数学研究所, 伊利诺伊大学芝加哥分校, 萨莱诺大学, 美因茨约翰内斯 - 古腾堡大学, 浦项科技大学以及安第斯大学讲授若干研究生短期课程, 讲义也随之不断扩充. 感谢参与者提供的宝贵反馈.

感谢 Ravi Vakil, Johan de Jong 和 Martin Olsson 引导我认识模空间与叠, 并鼓励我撰写这本教材. 若没有 de Jong 对 Stacks Project 的贡献, 以及他亲自与我分享的无数其他论证和例子, 本书便不可能完成; 同样, 若没有 Olsson 的教材 \cite{Ols16} 和他于 $2007$ 年在加州大学伯克利分校讲授的叠课程, 也不可能有本书, 这些工作塑造了我对这一学科的理解. 我也感谢合作者, 尤其是 Maksym Fedorchuk, Jack Hall, Daniel Halpern-Leistner, Jochen Heinloth, Andrew Kresch, David Rydh, David Smyth 和 Chenyang Xu, 感谢他们分享知识和见解. 感谢 Dhyan Aranha, Arend Bayer, Jim Bryan, Ben Church, Arkamouli Debnath, Hannah Dell, Neeraj Deshmukh, Patricio Gallardo, Davide Gori, Andres Fernandez Herrero, Vu Hien Giovanni Inchiostro, Adam Kapilow, Rega Karim, Myeonhu Kim, Quinn La Fond, Aaron Landsman, Dion Leijnse, Jake Levinson, Max Lieblich, Benjamin Lim, Siddarth Matthur, Andrés Ibáñez Núñez, Brian Nungent, Dragos Oprea, Lachlan Potter, Denis Resin, Minseon Shin, Tuomas Tajakka, Alex Schefflin, Ian Selvaggi, Reinier Sorgdrager, Dario Weissmann, Maxime Wybouw, Yujie Zhang, Xucheng Zhang 以及许多其他人提出的意见和更正.
